\documentclass[10pt]{article}

% graphicx package, useful for including eps, pdf, and raster graphics
\usepackage{graphicx}
\DeclareGraphicsExtensions{.png,.png,.jpg,.pdf}

% basic packages
\usepackage{fancyhdr}
\usepackage{color}
\usepackage{parskip}
\usepackage{float}
\usepackage{microtype}
\usepackage{url}
\usepackage{hyperref}
\usepackage{booktabs}
\usepackage{makecell}
\usepackage{multirow}
\usepackage{pbox}
\usepackage{enumitem}

% page layout
\usepackage[top=1in, bottom=1in, left=0.75in, right=0.75in, includefoot, heightrounded]{geometry}

% Useful for seeing the layout as you tweak it.
%\usepackage{showframe}

\setlength{\headheight}{1.1in}
\pagestyle{fancy}
\fancyhf{}

% fonts
\usepackage{fontspec}

\usepackage{polyglossia}
\setmainlanguage{russian}
\setotherlanguage{english}
\newfontfamily\cyrillicfont{DroidSans}[
  Extension   = .ttf,
  UprightFont = *,
  BoldFont    = *-Bold ]
\setmainfont{DroidSans}[
  Extension   = .ttf,
  UprightFont = *,
  BoldFont    = *-Bold ]
\setmonofont{DroidSansMono}[
  Extension   = .ttf,
  UprightFont = * ]

\hypersetup{colorlinks,urlcolor=blue}
\urlstyle{same}


\lhead{\includegraphics[height=1in,clip]{logo}}
\rhead{\pbox[b][1in][c]{\textwidth}{
    \small
    \textbf{Тест проводит следующая лаборатория:}\\
    Northwest Genomics Center\\
    University of Washington School of Medicine\\
    Genome Sciences Box 355065\\
    3720 15th Ave NE\\
    Seattle, WA 98195-5065}}

\newcommand{\PageLine}{\rule{\textwidth}{0.25mm}}

\newcommand{\link}[2]{\href{#1}{#2}\footnote{#1}}

\begin{document}

\begin{center}
\Large
\textbf{Отчёт лаборатории об исследовании Qual PCR NWGC SARS-CoV-2 (COVID-19) }
\end{center}

\bigskip

\begin{description}[font=\normalfont,align=left,labelwidth=18em]
\item [Ф.И.О. участника] \textbf{\VAR{pat_name|e}}
\item [Дата рождения] \textbf{\VAR{birth_date|e}}
\item [Идентификатор препарата] \textbf{\VAR{qrcode|e}}
\item [Дата предоставления образца] \textbf{\VAR{collect_ts|e}}
\item [Дата предоставления результатов]
  %- if status_code == "never-tested"
  \textbf{Обследование не проводилось}
  %- else
  \textbf{\VAR{result_ts|e}}
  %- endif
\item [Тип препарата] Средне-назальный, получен самостоятельно
\end{description}

\PageLine

Результат исследования Qual Polymerase Chain Reaction (PCR, полимеразной цепной
реакции) NWGC SARS-CoV-2 (COVID-19)

%- if status_code == "negative"
\textbf{Не обнаружен}\\
%- elif status_code == "positive"
\textbf{Обнаружен SARS-CoV-2 (положительный)}\\
%- elif status_code == "inconclusive"
\textbf{Неокончательный}\\
%- elif status_code == "never-tested"
\textbf{Не может тестироваться}\\
%- else
\VAR{0/0} %# poor man's raise
%- endif

\PageLine

Помните:

\begin{itemize}
\item

  \textbf{Сбор и тестирование образцов в данном проекте выполняются в рамках
  программы общественно-медицинского контроля, осуществляемой в партнёрстве с
  округом Кинг и другими учреждениями.} Полученный вами тестовый комплект не
  заменяет собой личной медицинской помощи со стороны квалифицированного
  работника здравоохранения.

\item

  \textbf{Ответственное хранение образца:} тестируемый образец был получен
  самостоятельно, самостоятельно промаркирован и отправлен в SCAN почтовой или
  курьерской службой. Его сбор, маркировка и перевозка выполнялись не
  лицензированным медицинским специалистом. Это может повлиять на результаты
  теста, причём нельзя исключать получения ложных отрицательных результатов.
  Может потребоваться дальнейшее тестирование в клинических условиях.

\item

  \textbf{Устройство для получения образца:} мазок из средней части носовой
  раковины не является типом образца, получившим Emergency Use Authorization
  (EUA, Разрешение для экстренного использования) от Food and Drug
  Administration (FDA, Управление по контролю качества пищевых продуктов и
  лекарственных препаратов)/Centers for Disease Control and Prevention (CDC,
  Центры по контролю и профилактике заболеваний). Поэтому то, что инфекция
  SARS-CoV-2 не обнаружена или получены отрицательные результаты, не может
  исключать наличия этой инфекции.

\item

  \textbf{«Positive» (Положительные) результаты могут быть подтверждены
  Washington State Department of Health (Департаментом здравоохранения штата
  Вашингтон).} Если у вас есть вопросы, рекомендуем обсудить данный результат с
  вашим врачом.

\item{
  \textbf{Если в результатах теста указано «Not detected» (Не обнаружено), это
  значит, что наша лаборатория не обнаружила SARS-CoV-2.}

  Все тесты несовершенны, поэтому существует небольшая вероятность того, что у
  вас всё же есть вирус SARS-CoV-2, который вызывает COVID-19 с отрицательным
  результатом.  При этом тесте не применяется стандартный способ сбора образца
  медицинским работником, и, кроме того, образец был отправлен почтой или
  курьерской службой, он не был получен и не перевозился в медицинских условиях.
  Независимо от результатов вашего теста, если вы почувствовали себя неважно,
  лучше оставаться дома и избегать контакта с другими людьми. Если у вас имеются
  другие нарушения здоровья, если вы серьёзно больны или беспокоитесь по поводу
  своей болезни, вам следует обратиться к своему врачу.  Если симптомы не
  проходят, обсудите со своим врачом возможность ещё раз пройти тест на
  COVID-19.
}

\item

  \textbf{Если в результатах теста указано «Inconclusive» (Неокончательный),}
  это означает, что полученные нашей лабораторией результаты теста на SARS-CoV-2
  носят неопределённый характер.

  Все тесты несовершенны, поэтому неопределённый результат может означать, что у
  вас всё же есть вирус SARS-CoV-2, который вызывает COVID-19.  При этом тесте
  не применяется стандартный способ сбора образца медицинским работником, и,
  кроме того, образец был отправлен почтой или курьерской службой, он не был
  получен и не перевозился в медицинских условиях.  Независимо от результатов
  вашего теста, если вы почувствовали себя неважно, лучше оставаться дома и
  избегать контакта с другими людьми. Если у вас имеются другие нарушения
  здоровья, если вы серьёзно больны или беспокоитесь по поводу своей болезни,
  вам следует обратиться к своему врачу.  Если симптомы не проходят, обсудите со
  своим врачом возможность ещё раз пройти тест на COVID-19.

\item

  \textbf{Если в результатах теста указано «Unable to be tested» (Не может
  тестироваться),} это означает, что либо образца было недостаточно, либо в
  процессе тестирования произошла ошибка.

  Поскольку тестирование не было выполнено, существует вероятность того, что у
  вас всё же есть вирус SARS-CoV-2, который вызывает COVID-19.  При этом тесте
  не применяется стандартный способ сбора образца медицинским работником, и,
  кроме того, образец был отправлен почтой или курьерской службой, он не был
  получен и не перевозился в медицинских условиях.  Независимо от результатов
  вашего теста, если вы почувствовали себя неважно, лучше оставаться дома и
  избегать контакта с другими людьми. Если у вас имеются другие нарушения
  здоровья, если вы серьёзно больны или беспокоитесь по поводу своей болезни,
  вам следует обратиться к своему врачу.  Если симптомы не проходят, обсудите со
  своим врачом возможность ещё раз пройти тест на COVID-19.

\item

  \textbf{О положительных результатах (с пометкой «DETECTED» (ОБНАРУЖЕНО) или
  «Virus Detected» (Обнаружен вирус)) исследовательское учреждение будет
  сообщать в соответствующие органы здравоохранения штата и округа.} Вам НЕ
  НУЖНО самим сообщать об этих результатах.  Если ваш тест указал на наличие у
  вас коронавируса нового типа, органы здравоохранения могут связаться с вами и
  задать вопросы о вашем заболевании и людях, с которыми вы могли вступать в
  контакт.

\item

  \textbf{Вы можете распечатать эти результаты теста или отправить их своему
  врачу по электронной почте.}

\end{itemize}

\bigskip
\textbf{Методика исследования Qual PCR SARS-CoV-2 (COVID-19):}

Цель этого теста –- обнаружить SARS-CoV-2 методом обратной транскрипции PCR. В
основе этого теста лежит усиление сегментов вирусного генома Orf1B и S с помощью
внутреннего контроля амплификации RNase P человека.  Данный тест был разработан
и его рабочие параметры были определены University of Washington’s Northwest
Genomics Center (Северо-Западным центром геномики Университета штата Вашингтон).
Данный тест, разработанный лабораторией, прошёл оценку и был представлен в US
Food and Drug Administration с целью получения Emergency Use Authorization. FDA
должно провести независимый анализ оценки теста. Данная лаборатория
сертифицирована в соответствии с Clinical Laboratory Improvement Amendments
(CLIA, Поправки к клиническим улучшениям в лабораторных условиях) и имеет право
на проведение клинических лабораторных тестов высокой степени сложности.

Чтобы получить дальнейшую информацию о медицинском учреждении перейдите по этим
ссылкам Washington State Department of Health:

\begin{itemize}
\item

  \link{https://www.doh.wa.gov/Portals/1/Documents/1600/coronavirus/COVIDcasepositiveRussian.pdf}{Пациенты
  с подтверждённым или подозреваемым случаем заболевания COVID-19}

\item

  \link{https://www.doh.wa.gov/Portals/1/Documents/1600/coronavirus/COVIDexposed\_notSickYet\_RU.pdf}{Пациенты,
  вступавшие в контакт с подтверждённым случаем заболевания COVID-19}

\item

  \link{https://www.doh.wa.gov/Portals/1/Documents/1600/coronavirus/COVIDsickNotExposedRussian.pdf}{Пациенты,
  не вступавшие в контакт, но имеющие симптомы COVID-19}

\end{itemize}

Дополнительные ресурсы опубликованы на веб-сайтах Washington State Department of
Health и CDC:

\begin{itemize}
\item

  \url{https://www.doh.wa.gov/Emergencies/Coronavirus/Russian}

\item

  \url{https://wwwn.cdc.gov/pubs/other-languages?Sort=Lang%3A%3Aasc&Language=Russian}

\end{itemize}

\end{document}
